\documentclass[conference]{IEEEtran}
\IEEEoverridecommandlockouts
\usepackage{cite}
\usepackage{amsmath,amssymb,amsfonts}
\usepackage{graphicx}
\usepackage{textcomp}
\usepackage{xcolor}
\usepackage{booktabs}
\usepackage{url}

% Numbers, p-values, sentences and tables produced by 03_inferential_stats.py
\input{generated/results_macros.tex}

\begin{document}

\title{A Statistical Analysis of Spatio-Temporal Traffic Travel Time Volatility in Dhaka Metropolitan Area}

\author{\IEEEauthorblockN{Shamiul}
\IEEEauthorblockA{\textit{Department of Computer Science} \\
Dhaka, Bangladesh \\
email address}
}

\maketitle

% =====================================================================
\begin{abstract}
Traffic congestion in Dhaka imposes large economic costs, yet quantitative evidence on how travel speed and its volatility differ across corridors and times of day remains sparse. This paper compares estimated travel speeds on \NCorr{} corridors of the Dhaka Metropolitan Area using \NObs{} snapshots logged every 15 minutes between \DateStart{} and \DateEnd{} by an automated OSRM and GitHub Actions pipeline. Speed is derived as distance divided by estimated duration, and corridor differences are tested with a one-way ANOVA supported by Shapiro--Wilk and Levene diagnostics, Welch's ANOVA, and Tukey HSD and Games--Howell post-hoc tests. Corridor accounted for a proportion $\eta^2 = \EtaSq$ of the variance in speed ($F(\AnovaDFB,\AnovaDFW) = \AnovaF$, $p \AnovaP$), and variances were \LevVerdict{} across corridors (Levene $W = \LevW$, $p \LevP$). \FastestCorr{} recorded the highest mean speed (\FastestMean~km/h) and \SlowestCorr{} the lowest (\SlowestMean~km/h). Travel-time volatility, measured by the coefficient of variation, was highest on \MostVolCorr{} (\MostVolCV\%) and lowest on \LeastVolCorr{} (\LeastVolCV\%). Across all corridors, mean speed was lowest in the \SlowSlot{} slot (\SlowSlotMean~km/h) and highest in the \FastSlot{} slot (\FastSlotMean~km/h). The open, reproducible pipeline provides a baseline for evidence-based traffic management in Dhaka.
\end{abstract}

\begin{IEEEkeywords}
traffic congestion, travel time variability, Dhaka, OSRM, one-way ANOVA, Tukey HSD, intelligent transportation systems
\end{IEEEkeywords}

% =====================================================================
\section{Introduction}
Dhaka is one of the densest and most congested metropolitan areas in the world. A World Bank assessment reports that average driving speed in the city fell from about 21~km/h to about 7~km/h over a decade and attributes part of the problem to a disorganised road network with no continuous east--west trunk road and no ring road \cite{worldbank2018dhaka}. A global comparison of more than 1{,}200 cities built from Google Maps travel times ranks Dhaka as the slowest city studied \cite{akbar2023fast}. Beyond the \emph{level} of congestion, commuters and planners are affected by its \emph{volatility}: the same trip can take very different times depending on the day and the hour, which undermines scheduling, logistics and public-transport reliability.

Most published Dhaka studies emphasise average speeds, cost estimates or network-level bottleneck detection \cite{sarker2025bottleneck}. Systematic, reproducible measurements that \emph{compare corridors} and \emph{describe how speed changes through the day} are less common, partly because commercial traffic feeds are expensive.

\subsection{Objectives}
The study has three objectives:
\begin{enumerate}
\item[O1] to build an automated, low-cost pipeline that logs estimated travel time and distance on selected Dhaka corridors every 15 minutes;
\item[O2] to test whether mean travel speed differs across corridors;
\item[O3] to quantify travel-time volatility on each corridor and describe how speed varies across four daily time slots (Morning Peak, Midday, Evening Peak, Night).
\end{enumerate}
The primary hypotheses are $H_0: \mu_1 = \mu_2 = \dots = \mu_k$ against $H_1$: at least one corridor mean speed $\mu_j$ differs, tested at $\alpha = 0.05$, where $k$ is the number of corridors.

\subsection{Contributions}
We contribute (i) an open pipeline based on OSRM \cite{luxen2011real} and GitHub Actions, (ii) an inferential analysis that pairs classical ANOVA with heteroscedasticity- and non-normality-robust alternatives, and (iii) a fully scripted workflow in which every table and statistic in this paper is generated from the data.

The remainder of the paper is organised as follows. Section~II reviews related work in South Asian cities, Section~III describes the methodology, Section~IV presents the results, and Section~V concludes.

% =====================================================================
\section{Literature Review}
\subsection{Dhaka, Bangladesh}
The World Bank's \emph{Toward Great Dhaka} report documents a decade-long collapse in average driving speed and links it to network structure and the slow progress of mass and bus rapid transit \cite{worldbank2018dhaka}. Akbar \emph{et al.} use travel times obtained from a mapping service at different times and days to compare speeds across more than 1{,}200 cities and find Dhaka to be the slowest, with Mymensingh and Chattogram also among the slowest \cite{akbar2023fast}. At a finer scale, Sarker \emph{et al.} identify bottlenecks in central Dhaka through graph-based congestion analysis and note that peak-hour speeds in parts of the city have fallen to 7--10~km/h \cite{sarker2025bottleneck}. These studies establish the level of congestion; the present study adds a formal statistical comparison across corridors and time slots.

\subsection{Indian metropolitan regions}
Fleet-based statistics for six Indian metropolitan regions in the first quarter of 2019 place average speeds between roughly 18 and 26~km/h, with Chennai the fastest and Mumbai and Bengaluru the slowest \cite{moveinsync2019}. Ghosh \emph{et al.} queried a mapping API for source--destination pairs in five Indian cities at 03:00, 09:00, 17:00 and 23:00 and found that, relative to 03:00, the extra time per kilometre at 09:00 was approximately 0.85 (Bengaluru), 0.74 (Mumbai), 0.70 (Delhi), 0.59 (Kolkata) and 0.57 (Hyderabad) \cite{ghosh2020eco}. This time-of-day contrast is the same quantity our time-slot analysis targets for Dhaka.

\subsection{Bengaluru corridor-level measurements}
A volunteer-led analysis of a 9~km Outer Ring Road segment in Bengaluru, based on live mapping-service traffic data collected for a week of morning and evening peaks, reported a free-flow time of 15 minutes but delays of up to 39 minutes on a Tuesday (about 10~km/h) versus up to 33 minutes on a Thursday \cite{citizenmatters}. This shows that large temporal differences are visible even on a single corridor, motivating a formal comparison across corridors and time slots.

\subsection{Methodological gap: static routing engines}
Ghosh \emph{et al.} also show that an OSRM-based engine with fixed speeds per road class produced durations equal to only about 0.58 (Bengaluru) and 0.53 (Delhi) of the mapping-service durations at 09:00, because the speed table does not vary with time of day \cite{ghosh2020eco}. This finding implies that OSRM-derived durations reflect time-varying congestion only if the routing back-end is supplied with time-dependent speed information, a limitation we return to in Section~V.

% =====================================================================
\section{Methodology}
\subsection{Study Corridors}
The study monitors \NCorr{} arterial corridors of the Dhaka Metropolitan Area listed in Table~\ref{tab:corridors}. Corridors are labelled C1, C2, \dots{} in decreasing order of mean speed, and the distance is the median route distance returned by the router.

\input{generated/corridor_table.tex}

\subsection{OSRM Data Collection Pipeline}
The Open Source Routing Machine (OSRM) computes shortest paths on a graph derived from OpenStreetMap \cite{luxen2011real}. A Python logger (\texttt{osrm\_logger.py}) requests a driving route for each corridor from the OSRM \texttt{route} service and records the returned distance and duration. The logger is executed by a GitHub Actions workflow on a 15-minute schedule and appends one row per corridor to a CSV file in the project repository. Each row contains \texttt{Timestamp}, \texttt{Day\_of\_Week}, \texttt{Corridor\_Name}, \texttt{Distance\_km} and \texttt{Estimated\_Duration\_Min}. Collection ran from \DateStart{} to \DateEnd{} (\NDays{} calendar days). Because GitHub schedules are best-effort, the number of snapshots per day is below the ideal 96 per corridor, so the number of observations differs between corridors (Table~\ref{tab:corridors}).

\subsection{Preprocessing}
The raw file is fetched from the repository and cleaned in the following order: (i) repeated header lines and empty or non-numeric API responses are removed; (ii) rows with non-positive duration or distance are removed; (iii) duplicate (\texttt{Timestamp}, \texttt{Corridor\_Name}) rows are removed; (iv) timestamps are converted to Asia/Dhaka local time (UTC+6) so that time slots are assigned by local hour; (v) rows whose distance deviates by more than 25\% from the corridor median are removed as route changes; and (vi) speeds above 120~km/h are removed as physically impossible. Extreme but plausible durations are retained, because volatility is the object of study. Speed in km/h is
\begin{equation}
v_i = \frac{d_i}{t_i/60},
\label{eq:speed}
\end{equation}
where $d_i$ is distance in km and $t_i$ is the estimated duration in minutes. Observations are assigned to the time slots in Table~\ref{tab:slotdef}, defined by local hour of day.

\begin{table}[!t]
\caption{Time-slot definitions (Asia/Dhaka local time)}
\label{tab:slotdef}
\centering\footnotesize
\begin{tabular}{ll}
\toprule
Time slot & Local time window \\
\midrule
Morning Peak & 07:00--09:59 \\
Midday & 10:00--15:59 \\
Evening Peak & 16:00--19:59 \\
Night & 20:00--06:59 \\
\bottomrule
\end{tabular}
\end{table}

\subsection{Statistical Design}
Let $y_{ij}$ be the speed of observation $i$ on corridor $j$ ($j = 1,\dots,k$), with $n_j$ observations per corridor, $N = \sum_j n_j$, corridor means $\bar{y}_j$ and grand mean $\bar{y}$. The one-way ANOVA statistic is
\begin{equation}
F = \frac{MS_B}{MS_W} = \frac{\sum_{j=1}^{k} n_j(\bar{y}_j - \bar{y})^2 / (k-1)}{\sum_{j=1}^{k}\sum_{i=1}^{n_j}(y_{ij} - \bar{y}_j)^2 / (N-k)},
\label{eq:anova}
\end{equation}
which follows an $F(k-1, N-k)$ distribution under $H_0$; the effect size is $\eta^2 = SS_B / SS_T$. Normality within each corridor is examined with the Shapiro--Wilk test \cite{shapiro1965analysis}, and homogeneity of variance with Levene's test using group medians \cite{levene1960robust}, which is also a direct test of whether speed volatility differs between corridors. Because both assumptions are strained by large, skewed samples, we also report Welch's ANOVA \cite{welch1951comparison} and the Kruskal--Wallis test \cite{kruskal1952use}. Pairwise differences use the Tukey--Kramer procedure \cite{tukey1949comparing} with studentised statistic
\begin{equation}
q = \frac{|\bar{y}_a - \bar{y}_b|}{\sqrt{\tfrac{MS_W}{2}\left(\tfrac{1}{n_a} + \tfrac{1}{n_b}\right)}},
\label{eq:tukey}
\end{equation}
and, since it does not assume equal variances, the Games--Howell procedure \cite{gameshowell1976pairwise}. Travel-time volatility on corridor $j$ is the coefficient of variation
\begin{equation}
\mathrm{CV}_j = 100 \times \frac{s_{t,j}}{\bar{t}_j},
\label{eq:cv}
\end{equation}
where $s_{t,j}$ and $\bar{t}_j$ are the standard deviation and mean of estimated travel time on corridor $j$. Successive 15-minute observations on a corridor are serially correlated, so the nominal $N$ overstates the information in the sample. We quantify this by the lag-1 autocorrelation $\rho_1$, approximate the effective sample size as $N_{\mathrm{eff}} \approx N(1-\rho_1)/(1+\rho_1)$, and repeat the ANOVA on daily corridor mean speeds as a robustness check.

\subsection{Software}
Analyses use Python with pandas \cite{mckinney2010data}, SciPy \cite{virtanen2020scipy} and statsmodels \cite{seabold2010statsmodels}; figures use matplotlib \cite{hunter2007matplotlib} and seaborn \cite{waskom2021seaborn}.

% =====================================================================
\section{Results}
\subsection{Descriptive Statistics}
The cleaned data contain \NObs{} observations. Table~\ref{tab:desc} and Fig.~\ref{fig:box} summarise speed by corridor. \FastestCorr{} had the highest mean speed (\FastestMean~km/h) and \SlowestCorr{} the lowest (\SlowestMean~km/h). Travel-time volatility ranged from $\mathrm{CV} = \LeastVolCV\%$ on \LeastVolCorr{} to $\mathrm{CV} = \MostVolCV\%$ on \MostVolCorr{}.

\input{generated/desc_table.tex}

\begin{figure}[!t]
\centering
\includegraphics[width=\columnwidth]{figures/fig1_speed_boxplot.pdf}
\caption{Distribution of travel speed by corridor (C1 = highest mean speed).}
\label{fig:box}
\end{figure}

\subsection{Assumption Checks}
The Shapiro--Wilk test rejected normality in \NormReject{} corridor groups at $\alpha = 0.05$; with samples of this size the test detects even trivial departures, so we rely on the robust tests below. Levene's test indicated \LevVerdict{} variances across corridors ($W = \LevW$, $p \LevP$). The lag-1 autocorrelation of speed within corridors was $\rho_1 = \RhoOne$, giving $N_{\mathrm{eff}} \approx \NEff$.

\subsection{One-Way ANOVA}
Table~\ref{tab:anova} reports the ANOVA. We \AnovaVerdict{} $H_0$: $F(\AnovaDFB,\AnovaDFW) = \AnovaF$, $p \AnovaP$, $\eta^2 = \EtaSq$. The robust alternatives give: Welch's ANOVA $F(\WelchDFnum,\WelchDFden) = \WelchF$, $p \WelchP$, and Kruskal--Wallis $H(\AnovaDFB) = \KWH$, $p \KWP$. Repeating the ANOVA on daily corridor mean speeds, which removes most of the serial dependence, gives $F(\AggDFB,\AggDFW) = \AggF$, $p \AggP$.

\input{generated/anova_table.tex}

\subsection{Post-Hoc Comparisons}
Of the \NPairs{} pairwise comparisons, \NSigTukey{} were significant under Tukey HSD and \NSigGH{} under Games--Howell. Table~\ref{tab:tukey} lists the pairs with the largest differences, and the complete set is available with the replication files. \SentTop{}

\input{generated/tukey_table.tex}

\subsection{Time-Slot Patterns}
Table~\ref{tab:slots} and Fig.~\ref{fig:slots} show mean speed by corridor and time slot. Pooled over corridors, mean speed was lowest in the \SlowSlot{} slot (\SlowSlotMean~km/h) and highest in the \FastSlot{} slot (\FastSlotMean~km/h).

\input{generated/slot_table.tex}

\begin{figure}[!t]
\centering
\includegraphics[width=\columnwidth]{figures/fig2_timeslot_lines.pdf}
\caption{Mean speed by time slot for each corridor. Error bars show 95\% bootstrap confidence intervals.}
\label{fig:slots}
\end{figure}

% =====================================================================
\section{Conclusion and Future Scope}
\subsection{Conclusion}
Using \NObs{} route snapshots from \NCorr{} Dhaka corridors, we tested whether mean speed differs across corridors. The omnibus ANOVA led us to \AnovaVerdict{} the null hypothesis ($F(\AnovaDFB,\AnovaDFW) = \AnovaF$, $p \AnovaP$, $\eta^2 = \EtaSq$), with \FastestCorr{} fastest and \SlowestCorr{} slowest on average. Volatility of travel time, measured by CV, was highest on \MostVolCorr{} and lowest on \LeastVolCorr{}, and the time-slot profile in Fig.~\ref{fig:slots} documents how speed changes through the day.

\subsection{Limitations}
First, the OSRM default car profile assigns speeds by road class rather than from live traffic \cite{ghosh2020eco}; the logged durations are therefore model-based estimates, not measured travel times, and any temporal variation reflects only the time-dependent information available to the routing back-end. Second, observations are serially correlated, so nominal $p$-values are optimistic and effect sizes should be emphasised. Third, the collection window is short (\NDays{} calendar days), so date-specific events such as weather and public holidays may influence the corridor comparison. Fourth, corridors differ in length and road class, so differences in mean speed describe the corridors themselves and not a causal effect of any single factor.

\subsection{Future Scope}
Future work will (i) validate estimates against a traffic-aware provider or GPS probe data, (ii) replace one-way ANOVA with mixed-effects models with corridor random intercepts and autoregressive errors, (iii) extend collection over Ramadan, Eid and monsoon periods, (iv) add covariates such as day of week, rainfall and public events, and (v) study spatial spill-over between adjacent corridors.

% =====================================================================
\bibliographystyle{IEEEtran}
\bibliography{references}

\end{document}
